\section{引言}
\label{sec:1}

组合——从简单部件组装复杂系统——是软件工程的一个基本原则\cite{ref1}。传统上，组合是静态的：函数调用、模块导入和类继承在编译时解析，并在整个执行过程中保持固定。然而，现代软件日益需要动态组合，即组件在运行时被加载、卸载和重新配置。插件架构\cite{ref2}和自演化智能体框架都需要能够在运行中安全地添加和移除功能的系统，但当前实践依赖粗粒度机制\cite{ref3}，只能通过重启来重新配置，丢弃运行时状态。尽管动态组合的实际重要性日益增长，与静态组合已有的丰富形式化框架相比，其理论基础仍然不成熟。

\subsection{可组合性的维度}
\label{sec:1.1}

为了刻画动态组合的需求，我们在已被充分研究的组合代数层面之外，识别出两个正交的维度：

\begin{itemize}
\item 时间可组合性处理时间维度：在移除组件时，该组件对共享环境所做的修改必须被完全且安全地撤销。这要求追踪组件执行的每一项资源分配、事件注册和状态变更，并保证在移除时有序地回收它们。
\item 空间可组合性处理空间维度：组件必须能够以结构化且可验证的方式声明、发现和解析彼此之间的依赖关系。这要求管理依赖拓扑，并协调组件生命周期以响应依赖变化。
\end{itemize}

在静态场景中，时间可组合性归约为词法作用域（例如 RAII\cite{ref4}、括号模式\cite{ref5}），空间可组合性归约为模块导入解析\cite{ref6}。在动态场景中，组件在运行时到达和离去，两个维度都变得显著更加困难：时间可组合性必须处理作用域不受词法约束的长期有状态效应；空间可组合性必须处理在执行期间出现、消失或改变身份的依赖关系。

\subsection{动机示例}
\label{sec:1.2}

\subsubsection{插件系统}
\label{sec:1.2.1}

插件系统是动态组合的典型实例。我们以 Visual Studio Code（VSCode）——最广泛使用的可扩展 IDE 之一——作为代表性示例。

时间局限性。VSCode 在一个名为扩展宿主的共享进程中运行所有扩展。虽然扩展可以动态安装，但该宿主没有提供在运行时卸载单个扩展代码的机制。一旦扩展的 activate 函数被执行，禁用或卸载它就需要重启整个宿主，影响所有已加载的扩展。纯粹声明性的扩展（如主题、快捷键绑定和代码片段）不携带代码，可以自由移除。然而，在安装数量排名前 100 的扩展中，87 个包含可执行代码\footnote{数据于2026年6月9日从 Visual Studio Code Marketplace 获取。}，因此在移除时需要这样的重启。尽管 VSCode 提供了 deactivate 钩子，它仅作为宿主进程终止期间的优雅关闭回调，因此不支持实时移除。此外，该钩子将效应处置与效应创建（在 activate 中）分离，违反了关注点局部性，使得完整清理难以验证。

空间局限性。VSCode 确实提供了 extensionDependencies 用于声明扩展之间的依赖关系，但使用极少：在安装数量排名前 100 的扩展中，仅 7 个声明了对非内置扩展的 extensionDependencies\footnote{数据于2026年6月9日从 Visual Studio Code Marketplace 获取。}。这种稀少反映了扩展 API 的形态，它暴露固定的、表层的扩展点，如命令、视图和语言特性。扩展通过这些点向宿主贡献功能，而非彼此依赖，因此扩展间的依赖很少出现。此外，VSCode 的扩展间交互机制不提供结构化契约：它通过 vscode.extensions.getExtension(...).exports 将扩展的功能暴露给其他扩展，但返回值是无类型的（默认为 any），因此依赖方无法依赖经过检查的接口。简言之，VSCode 将扩展导向一组固定的宿主提供的扩展点，并未为它们彼此依赖提供安全、结构化的方式。

这两个局限性并非 VSCode 所独有；它们在一般插件系统中反复出现\cite{ref2,ref7}，仅在程度上有所不同。

\subsubsection{自演化智能体框架}
\label{sec:1.2.2}

现代 AI 智能体依赖运行时智能体框架\cite{ref8,ref9,ref10}。这些系统可能组合多样的工具套件\cite{ref11}和执行环境，管理权限和沙箱，维护会话状态和持久性，提供上下文管理和记忆系统\cite{ref12}，编排子智能体和多智能体工作流\cite{ref13}，并向用户和自动化暴露接口。未来的框架可能在持续服务请求的同时，生成并部署对其自身组件的修改。模型合成的可复用工具为组件级自修改提供了一个更窄的前身\cite{ref14}。每一次这样的修改本身就是动态组合的一个实例。

由于这些修改持续发生且人工监督有限或没有，动态可组合性变得不可或缺。没有时间可组合性，每次自修改都迫使完全重启，丢弃所有进程本地累积的状态；在这种频率下，累积的不可用时间变得可观，且进行中的任务被反复中断；更糟的是，一个有缺陷的自修改可能禁用恢复所需的过程本身。没有空间可组合性，每个模块必须自行检测并适应其所依赖模块在出现、消失或改变身份时的变化，且只能以临时方式实现；更糟的是，一种朴素的代码替换策略可能静默地破坏依赖方，或引入仅在重载时才显现的循环依赖。

\subsubsection{粗粒度权宜之计}
\label{sec:1.2.3}

动态可组合性受到有限形式化关注的一个原因是，操作系统和容器编排器已经提供了粗粒度的替代。操作系统以进程的粒度提供时间可组合性；容器编排器\cite{ref3}以服务的粒度提供空间可组合性。在实践中，大多数软件通过依赖这些粗粒度机制来容忍细粒度可组合性的缺失：行为异常的模块通过重启进程来处理，服务依赖由容器编排器管理。

然而，这种权宜之计带来了巨大的代价。在时间上，每次重启丢弃所有进程本地累积的状态（例如缓存、连接、部分计算），重建需要数秒到数分钟\cite{ref15}；在此期间维持可用性需要冗余副本，产生资源开销以补偿无法恢复单个组件的缺陷。在空间上，容器级编排无法表达共享地址空间的组件之间的依赖关系，并为本可以是本地函数调用的交互引入了网络开销。两种机制都在进程和容器的边界上运作，但现代系统日益在更细的层面上进行组合。这种粒度失配需要一种组合抽象，在与组件自身相同的层面上管理效应和依赖。

\subsection{贡献}
\label{sec:1.3}

动态可组合性的两个维度分别关注计算如何修改其环境以及如何依赖于其环境。这两个方向正是效应系统\cite{ref16,ref17}和协效应系统\cite{ref18,ref19}所形式化的内容：效应提供了关于环境修改推理的形式化词汇，协效应则提供了关于环境需求推理的形式化词汇。然而，现有的形式化将推理限制在词法固定作用域上的编译时分析，并未扩展到组件在运行时到达和离去的动态场景。通过将效应提升为可逆的运行时模型、将协效应提升为反应式依赖解析机制，我们获得了动态可组合性的统一形式化基础，该基础与语言无关，适用于任何需要动态组合的软件架构。我们做出以下贡献：

\begin{enumerate}
\item 我们形式化了可逆效应（\S\ref{subsec:3.1}）：每个上下文变换携带一个运行时追踪的显式逆操作，且追踪和恢复均保持组合性，因此在组件移除时上下文得以恢复。这建立了局部时间可组合性。
\item 我们形式化了反应式协效应（\S\ref{subsec:3.2}）：组件将其所需的协效应声明为规约，上下文的每次变更根据该规约通知组件为激活、去激活或中性的。这建立了局部空间可组合性。
\item 我们将效应上下文和协效应上下文统一为单一的上下文类型（\S\ref{subsec:3.3}），其中协效应上的观察等价性为效应提供了独立性，构成了时空可组合性的编程范式。
\item 我们给出了动态组合的演算（\S\ref{sec:4}），将两种机制组合为组件的概念，并为其生命周期配备了操作语义。其元理论将时空可组合性从单个组件推广到交织组件的整个系统。
\item 我们在 Cordis（\S\ref{sec:5}）中实现了这些思想，这是一个时空可组合性的元框架，提供了实现形式化模型的核心库（包含效应追踪和协效应解析），以及具有配置调和和热模块替换的声明式组件加载器。
\end{enumerate}
